\section{Component labels and annular depth}
\label{sec:annuli}

The graph paths supplied by Proposition~\ref{block:setting} stay near a
boundary point, but they may pass through vertices that we need to remove.
We first show how to assign a component of the remaining graph to a nearby
boundary point. We then construct an invariant integer-valued depth on
group vertices. These are the two inputs for pushing paths towards the
boundary in the next section.

We use the following setting from Proposition~\ref{block:setting}.
The compact metric space $(Y,\rho)$ is the disjoint union $H\sqcup L$,
where $H$ is a group, its points are isolated, and $H$ acts on $Y$ by
homeomorphisms extending left multiplication. The space $L$ is a
nondegenerate continuum, has no cut point, and satisfies
$\dim L\leq 1$ and $\check H^1(L;\mathbb F_2)=0$.
If distinct elements $g_n\in H$ satisfy $g_n\to a$ and
$g_n^{-1}\to b$, then $a,b\in L$ and $g_n$ converges to $a$ uniformly
on compact subsets of $Y\setminus\{b\}$.

There is a connected $H$-graph $K$ with vertex set $I$, where either
$I=H$ or $I\subset L$. It has finitely many vertex and edge orbits.
Every point of $L$ is a limit of distinct vertices of $K$, and, for every
$x\in L$ and neighborhood $W$ of $x$ in $Y$, there is a neighborhood
$V$ of $x$ such that any two vertices in $I\cap V$ can be joined by a
finite path of $K$ whose vertices lie in $W$.
For each fixed $w\in I$,
\[
  \rho(gw,g)\longrightarrow 0
  \qquad(g\in H\text{ leaving every finite set}).
\]
When $I\subset L$, every $p\in I$ has an infinite-order stabilizer
element $h\in H$ for which $h^jz\to p$ for every $z\in L$ as
$j\to+\infty$. No local finiteness of $K$ or finite generation of $H$
is assumed here.

\subsection{Extending labels across a deleted point}

We write $k=\mathbb F_2$ with its discrete topology. A locally constant
map to $k$ records a division into two open and closed parts.
The cohomology hypothesis prevents such a division from appearing only
after one point has been removed from a small neighborhood.
The circle shows why absence of cut points alone would not suffice:
removing the midpoint of a small open arc leaves two sides that can
carry different constant labels.

\begin{lemma}\label{ann:extension}
Let $L$ be a nondegenerate compact metric continuum with no cut point and
$\check H^1(L;k)=0$. If $x\in L$ and $V$ is an open neighborhood of $x$
in $L$, every locally constant map $u:V\setminus\{x\}\to k$ extends
uniquely to a locally constant map on $V$.
\end{lemma}

\begin{proof}
Put $U=L\setminus\{x\}$. We prove directly, using finite-cover
cohomology, that there are locally constant maps $f_U:U\to k$ and
$f_V:V\to k$ such that
\[
  u=f_U+f_V\quad\text{on }U\cap V.
\]
Choose open sets $U_0,V_0$ covering $L$ whose closures are contained in
$U,V$, respectively. The set
$C=\overline{U_0}\cap\overline{V_0}$ is a compact subset of $U\cap V$.
Its two subsets $C\cap u^{-1}(0)$ and $C\cap u^{-1}(1)$ are disjoint
compact sets. Choose $\delta>0$ smaller than their distance when both
are nonempty; if either is empty, any positive $\delta$ suffices.

Take a finite open cover $(W_i)$ of $L$ refining $(U_0,V_0)$, with each
$W_i$ of diameter less than $\delta$, and give it a parent
$P_i\in\{U,V\}$ according to this refinement. On every nonempty
$W_i\cap W_j$, define $c_{ij}=0$ if $P_i=P_j$, and define $c_{ij}$ to
be the value of $u$ there otherwise. This latter value is constant:
the intersection lies in $C$ and cannot meet both of its two compact
parts. At a triple intersection either all parents agree, or exactly
two of the three transition values equal the same value of $u$.
Thus $(c_{ij})$ is a $1$-cocycle on the nerve of $(W_i)$.

Since $\check H^1(L;k)=0$, some finite open refinement $(Z_\alpha)$,
with refinement map $r$, admits constants $a_\alpha\in k$ satisfying
\[
  a_\alpha+a_\beta=c_{r(\alpha)r(\beta)}
  \quad\text{whenever }Z_\alpha\cap Z_\beta\ne\varnothing.
\]
For $z\in U\cap Z_\alpha$, set $f_U(z)=a_\alpha$ if
$P_{r(\alpha)}=U$, and set $f_U(z)=a_\alpha+u(z)$ otherwise.
The displayed identity makes these definitions agree on overlaps.
They are locally constant. Define $f_V$ by the same rule, with $U,V$
interchanged. Then $f_U+f_V=u$ on $U\cap V$.

Since $x$ is not a cut point, $U$ is connected, so $f_U$ has one
constant value, say $c$. The map $f_V+c$ is the required extension.
Uniqueness follows because a nondegenerate continuum has no isolated
points, so $V\setminus\{x\}$ is dense in $V$.
\end{proof}

For a set $D\subset Y$, denote by $D'$ the set of limits in $Y$ of
sequences of distinct points of $D$. In a metric space $D'$ is closed.
For an allowed vertex set $T\subset I$, put
\[
  \Omega_T=L\setminus(I\setminus T)'.
\]
This is precisely the open set of boundary points near which deleted
vertices do not accumulate as distinct points. We use graph components
of the induced graph $K[T]$.

\begin{lemma}\label{ann:component}
In the setting above, each $x\in\Omega_T$ determines a unique component
$Q_T(x)$ of $K[T]$ such that all allowed vertices in some neighborhood
of $x$ belong to $Q_T(x)$. Every $x\in\Omega_T$ is a limit of distinct
allowed vertices, and the assignment $Q_T$ is locally constant on
$\Omega_T$. If $t\in T\cap L\cap\Omega_T$, then $Q_T(t)$ is the
component containing the vertex $t$ itself.
\end{lemma}

\begin{proof}
There is a neighborhood $W$ of $x$ such that
$(I\setminus T)\cap W\subset\{x\}$. Since $x$ is a limit of distinct
vertices of $I$, it is consequently a limit of distinct vertices of
$T$ as well.

Suppose first that $x$ is not a deleted vertex. After shrinking $W$,
it contains no deleted vertex. The near-path property then puts all
vertices sufficiently near $x$ in one component of $K[T]$. There are
allowed vertices arbitrarily near $x$, so this component is unique.
The same neighborhood, together with approximation by allowed vertices,
shows that the component assignment is locally constant near $x$.
When $x$ itself is an allowed vertex, it belongs to this component.

The remaining case has $I\subset L$ and $x\in I\setminus T$.
Choose an open neighborhood $V$ of $x$ in $L$ containing no other
deleted vertex. By the preceding case, every $y\in V\setminus\{x\}$
has a component label, and these labels form a locally constant map
\[
  q:V\setminus\{x\}\longrightarrow\{\text{components of }K[T]\}.
\]
Every map from this label set to $k$, composed with $q$, extends across
$x$ by Lemma~\ref{ann:extension}. We claim this forces $q$ itself to
be constant sufficiently near $x$.

To see this, call a label recurrent if it occurs arbitrarily near $x$.
If two labels are recurrent, a coloring separating them contradicts
the extension property. If exactly one label is recurrent but $q$ is
not eventually constant, color that label $0$ and all others $1$ to
obtain the same contradiction. If no label is recurrent, a sequence
of points tending to $x$ can be chosen with pairwise distinct labels:
at each step, first exclude the finitely many labels already chosen.
Color these successive labels alternately $0$ and $1$.
This again contradicts extension as a locally constant map.
Thus one component label occurs throughout a small punctured
neighborhood of $x$. Every allowed vertex there belongs to its own
label by the first case, so all these vertices belong to that component.

This proves existence at a deleted vertex too. Approximation by allowed
vertices proves uniqueness and, on shrinking the neighborhood once
more, local constancy of $Q_T$ there.
\end{proof}

\subsection{An invariant system of annuli}

We use the annulus-system method of
\cite[Sections~6--7]{Bowditch1998Topological}, with the following
nesting convention. The construction and the two properties needed
here are proved explicitly; the sides need not be connected.

An \emph{annulus} in $Y$ is an ordered pair
$\sigma=(\sigma^-,\sigma^+)$ of disjoint closed subsets of $Y$ such
that each side meets $L$ and $\sigma^-\cup\sigma^+\ne Y$. Its sides
need not be connected. For annuli $\sigma,\tau$, write
\[
  \sigma<\tau\quad\Longleftrightarrow\quad
  \sigma^+\cup\tau^-=Y.
\]
Then $\sigma^-\subset\operatorname{int}(\tau^-)$ and
$\tau^+\subset\operatorname{int}(\sigma^+)$. The relation is transitive
and irreflexive. Thus a nested chain
$\sigma_1<\cdots<\sigma_n$ contains distinct annuli.
A chain goes \emph{from $F$ to $E$} if
$F\subset\sigma_1^-$ and $E\subset\sigma_n^+$.

\begin{lemma}\label{ann:system}
There is a set $\mathcal A$ of annuli invariant under $H$ and under
interchange of the two sides, with the following properties.
\begin{enumerate}
\item For every $\delta>0$, only finitely many annuli in $\mathcal A$
have both sides of diameter at least $\delta$.
\item For every closed $F\subset Y$, every $x\in L\setminus F$, and
every integer $n\geq1$, there is a chain of $n$ annuli in $\mathcal A$
from $F$ to a neighborhood of $x$ in $Y$.
\end{enumerate}
\end{lemma}

\begin{proof}
We first record a uniform consequence of convergence. For every
$r,\varepsilon>0$, all but finitely many $g\in H$ admit a point
$b_g\in L$ such that
\[
  \operatorname{diam}g\bigl(Y\setminus B_\rho(b_g,r)\bigr)
  <\varepsilon.
\]
Otherwise choose distinct exceptions $g_n$ and pass to a subsequence
with $g_n\to a\in L$ and $g_n^{-1}\to b\in L$. Uniform convergence
on the compact set $Y\setminus B_\rho(b,r)$ gives a contradiction.

Choose $r_i\downarrow0$ with $r_i<\operatorname{diam}(L)/3$, and put
$\varepsilon_i=2^{-i}$. Apply the preceding observation with
$r=r_i/4$ and $\varepsilon=\varepsilon_i$, obtaining a finite
exceptional set $E_i\subset H$. Uniform continuity of the finitely
many maps in $E_i$ permits a common choice $0<s_i<r_i/4$ such that
\[
  \operatorname{diam}g\bigl(\overline{B_\rho(y,s_i)}\bigr)
  <\varepsilon_i
  \quad(g\in E_i,\ y\in L).
\]
Choose finitely many centers $y_{ij}\in L$ whose $s_i/2$-balls cover
$L$. At each center take the seed annulus
\[
  \sigma_{ij}^- =Y\setminus B_\rho(y_{ij},r_i),
  \qquad
  \sigma_{ij}^+ =\overline{B_\rho(y_{ij},s_i)}.
\]
Its sides are disjoint and have distance at least $r_i-s_i$.
The positive side contains its center. The negative side meets $L$
because every point of $L$ is at distance at least
$\operatorname{diam}(L)/2$ from some point of $L$.
The sides cannot cover $L$, since their nonempty closed intersections
with $L$ would disconnect it.

Let $\mathcal A$ consist of all distinct translates of these annuli
and their reversals. A ball of radius $r_i/4$ cannot meet both sides
of a seed annulus, whose mutual distance exceeds $3r_i/4$.
For $g\notin E_i$, one side therefore lies outside the exceptional
ball for $g$ and has image of diameter less than $\varepsilon_i$.
For $g\in E_i$, the positive side has that property by the choice of
$s_i$. Consequently every translated level-$i$ annulus has at least
one side of diameter less than $\varepsilon_i$.

For a fixed seed and fixed $\delta>0$, only finitely many distinct
translates can have both sides of diameter at least $\delta$.
If not, choose corresponding distinct group elements and a convergence
subsequence with repeller $b$. At least one of the two closed seed
sides excludes $b$, so its images have diameter tending to zero.
There are only finitely many seeds at each level, and levels with
$\varepsilon_i<\delta$ contribute no annuli with two such large sides.
This proves the first assertion.

For the second assertion, let $F$ be closed and $x\in L\setminus F$.
At a sufficiently fine level, choose a center $y_{ij}$ with
$\rho(x,y_{ij})<s_i/2$ and
$r_i+s_i/2<\operatorname{dist}(x,F)$ when $F$ is nonempty.
The negative side contains $F$, and the positive side contains a
neighborhood of $x$. Suppose an annulus already chosen has positive
side containing $B_\rho(x,\eta)$. Choose a finer level with
$r_j+s_j/2<\eta$, and a center $y$ within $s_j/2$ of $x$.
Then $B_\rho(y,r_j)\subset B_\rho(x,\eta)$, so the new annulus is
larger in the order $<$, and its positive side again contains a
neighborhood of $x$. Repeating gives any prescribed finite length.
\end{proof}

Fix such a system $\mathcal A$. For $F\subset Y$ and $v\in Y$, define
\[
  D(F,v)=\sup\{n\geq1:
       \text{a chain of length }n\text{ in }\mathcal A
       \text{ goes from }F\text{ to }\{v\}\},
\]
with value $0$ if there is no such chain. In general the value belongs
to $\{0,1,2,\ldots\}\cup\{\infty\}$. We abbreviate $D(\{u\},v)$
to $D(u,v)$. This function counts separating annuli; it is not asserted
to be a metric.

\begin{lemma}\label{ann:depth}
For $g\in H$, $F\subset Y$, and $v\in Y$,
\[
  D(gF,gv)=D(F,v).
\]
If $F$ contains a point of $H$ and $v\in H$, then $D(F,v)$ is finite.
When it is positive, a chain attains this finite maximum.
If $F$ is closed, $x\in L\setminus F$, and $N\geq1$, there is a
neighborhood $V_N$ of $x$ in $Y$ such that $D(F,v)\geq N$ for every
$v\in V_N$.
\end{lemma}

\begin{proof}
Translation carries annuli and chains bijectively to annuli and chains,
which proves equivariance. Suppose $h\in F\cap H$ and $v\in H$.
As $L$ is closed and excludes $h,v$, the number
\[
  \delta=\min\{\operatorname{dist}(h,L),
                    \operatorname{dist}(v,L)\}
\]
is positive. Throughout a chain from $F$ to $v$, every negative side
contains $h$ and meets $L$, and every positive side contains $v$ and
meets $L$. Both sides therefore have diameter at least $\delta$.
Lemma~\ref{ann:system} supplies only finitely many such annuli, and a
strict chain cannot repeat one. This proves finiteness. A finite
nonzero supremum of possible integer lengths is attained.
The final assertion is precisely the second part of
Lemma~\ref{ann:system}, using its terminal neighborhood.
\end{proof}
